\documentclass[11pt, a4paper]{article}

\usepackage[utf8]{inputenc}
\usepackage[T1]{fontenc}
\usepackage{lmodern}
\usepackage{amsmath, amssymb, amsthm, mathtools}
\usepackage[margin=1in]{geometry}
\usepackage{enumitem}
\usepackage{hyperref}
\usepackage{xcolor}
\usepackage{tikz}
\usepackage{tikz-cd}
\usepackage{booktabs}
\usepackage{longtable}
\usepackage{microtype}
\usepackage{listings}
\usepackage{fancyvrb}

\hypersetup{
    colorlinks=true,
    linkcolor=blue!50!black,
    citecolor=green!50!black,
    urlcolor=blue!50!black,
    pdftitle={Verum as a Foundation-Neutral Host of the Moduli Space of Formal Systems},
    pdfauthor={Max Sereda}
}

\tolerance=1500
\emergencystretch=2em
\hyphenpenalty=200

% Theorem environments
\newtheorem{theorem}{Theorem}[section]
\newtheorem{proposition}[theorem]{Proposition}
\newtheorem{lemma}[theorem]{Lemma}
\newtheorem{corollary}[theorem]{Corollary}
\theoremstyle{definition}
\newtheorem{definition}[theorem]{Definition}
\newtheorem{example}[theorem]{Example}
\newtheorem{construction}[theorem]{Construction}
\theoremstyle{remark}
\newtheorem{remark}[theorem]{Remark}
\newtheorem{convention}[theorem]{Convention}

% MSFS macros
\newcommand{\RS}{\mathrm{R\text{-}S}}
\newcommand{\cS}{\mathcal{S}}
\newcommand{\cF}{\mathcal{F}}
\newcommand{\cM}{\mathcal{M}}
\newcommand{\cC}{\mathcal{C}}
\newcommand{\cD}{\mathcal{D}}
\newcommand{\cL}{\mathcal{L}}
\newcommand{\cP}{\mathcal{P}}
\newcommand{\fM}{\mathfrak{M}}
\newcommand{\Meta}{\mathfrak{Meta}}
\newcommand{\LFnd}{\mathcal{L}_{\mathrm{Fnd}}}
\newcommand{\LCls}{\mathcal{L}_{\mathrm{Cls}}}
\newcommand{\LClsMax}{\mathcal{L}_{\mathrm{Cls}}^{\top}}
\newcommand{\LAbs}{\mathcal{L}_{\mathrm{Abs}}}
\newcommand{\Sint}{\cS_{\mathrm{int}}}
\newcommand{\Cl}{\mathrm{Cl}}
\newcommand{\ZFC}{\mathsf{ZFC}}
\newcommand{\MLTT}{\mathsf{MLTT}}
\newcommand{\HoTT}{\mathsf{HoTT}}
\newcommand{\ETT}{\mathsf{ETT}}
\newcommand{\CIC}{\mathsf{CIC}}
\newcommand{\II}{\mathbf{I}}
\newcommand{\Eff}{\mathrm{Eff}}
\newcommand{\Verum}{\textsc{Verum}}
\newcommand{\Coq}{\textsc{Coq}}
\newcommand{\Lean}{\textsc{Lean}}
\newcommand{\Agda}{\textsc{Agda}}
\newcommand{\Isabelle}{\textsc{Isabelle}}
\newcommand{\Fstar}{\textsc{F$^\ast$}}
\newcommand{\Pellet}{\textsc{Pellet}}
\newcommand{\HermiT}{\textsc{HermiT}}
\newcommand{\FactPP}{\textsc{Fact++}}
\newcommand{\ELK}{\textsc{ELK}}
\newcommand{\Konclude}{\textsc{Konclude}}
\newcommand{\OWL}{\textsc{OWL\,2}}
\newcommand{\SUMO}{\textsc{SUMO}}
\newcommand{\Cyc}{\textsc{Cyc}}
\newcommand{\DOLCE}{\textsc{DOLCE}}
\newcommand{\BFO}{\textsc{BFO}}
\newcommand{\SWRL}{\textsc{SWRL}}
\newcommand{\RDF}{\textsc{RDF}}

% Code listing style
\lstdefinelanguage{verum}{
    keywords={fn, let, is, type, protocol, implement, if, else, match, for, while, return, mount, module, public, mut, Self, self, true, false, in, using, where, requires, ensures, async, await, axiom, theorem, lemma},
    keywordstyle=\color{blue!70!black}\bfseries,
    ndkeywords={Int, Bool, Float, Text, List, Map, Set, Maybe, Result, Heap, Shared, Path, Equiv},
    ndkeywordstyle=\color{purple!60!black},
    sensitive=true,
    comment=[l]{//},
    morecomment=[s]{/*}{*/},
    commentstyle=\color{gray!70!black}\itshape,
    stringstyle=\color{green!40!black},
    string=[b]",
    morestring=[b]',
    otherkeywords={@verify, @framework, @derive, @pure, @inline, @const}
}
\lstset{
    language=verum,
    basicstyle=\ttfamily\small,
    breaklines=true,
    showstringspaces=false,
    frame=single,
    framesep=4pt,
    rulecolor=\color{gray!30},
    captionpos=b,
    numbers=none,
}

\title{%
    \LARGE \textbf{\Verum{} as a Foundation-Neutral Host of the Moduli Space of Formal Systems} \\[4pt]
    \large Compiler-Integrated Verification Architected under AFN-T
}
\author{Max Sereda \\ \small Independent researcher}
\date{2026}

\begin{document}

\maketitle

\begin{abstract}
\noindent
The proof assistant \Verum{} is analysed as a \emph{foundation-neutral host} for the moduli space $\fM$ of Rich-foundations classified by the MSFS programme. Unlike every proof assistant to date, \Verum{} does not commit to a single foundation; instead, its standard library provides a domain-neutral classifier apparatus (\texttt{core.theory\_interop}, \texttt{core.verify}, \texttt{core.proof}, \texttt{core.math.frameworks}) that accepts any $\LFnd$-point as an external input. The result is four structural correspondences proved here: (i) the \texttt{@framework(name, citation)} mechanism realises an indexing functor into $\LFnd$; (ii) \texttt{protocol}-based abstraction realises the classifier stratum $\LCls$ with explicit (M1)--(M5) data; (iii) the LCF-style trusted kernel embodies T-2f* depth-stratification and inherits universal paradox-immunity via Yanofsky-reducibility; (iv) the nine-strategy \texttt{@verify(...)} dispatch refines the Diakrisis $\nu$-invariant in an operational-cost gradient. The correspondences are demonstrably domain-neutral: neither the language nor its stdlib privileges any specific application, and any external cog satisfying the corpus-pattern automatically inhabits its own MSFS coordinate. The paper extends this framework to the knowledge-representation landscape: \OWL-\textsc{DL}, \SWRL, modal / epistemic / temporal / dynamic / $\mu$-calculus, default / circumscription / answer-set, probabilistic / fuzzy / possibilistic are all hosted as external \texttt{@framework}-packages with computable MSFS coordinates; industrial tableau reasoners (\Pellet, \HermiT, \FactPP, \ELK, \Konclude) are subsumed on their decidable fragments with kernel-verified certification while cross-paradigm translation via \texttt{core.theory\_interop} gauge-transformations remains structurally unavailable elsewhere. Beyond integration, the paper argues that \Verum{} occupies a position no extant proof assistant or reasoner can match: \Coq, \Lean, \Agda, \Isabelle, \Fstar, Dafny, Liquid Haskell, Mizar, Metamath each entangle language and foundation; OWL2 reasoners commit to SROIQ alone; \Verum{} separates language from foundation and hosts the full KR ladder. The gap is architectural orthodoxy under MSFS, not engineering debt.
\end{abstract}

\tableofcontents

\section{Introduction}
\label{sec:intro}

Proof assistants have accumulated for six decades. Each embodies a specific foundation: \Coq{} / \textsc{Rocq} realises the Calculus of Inductive Constructions; \Lean{} 4 realises a dependent type theory with proof irrelevance; \Agda{} ships cubical type theory natively; \Isabelle{} / HOL implements a weakly classical higher-order logic; \Fstar{} couples refinement types with monadic effects; Dafny targets SMT-discharged contracts; Liquid Haskell decorates Haskell with refinement predicates; Mizar formalises Tarski--Grothendieck set theory; Metamath hosts a minimal substitution calculus over ZFC.

Every one of these systems \emph{is} a foundation. The user who imports \Coq{} imports CIC; the user who imports \Lean{} imports Lean's dependent type theory; no separation is available. This architectural choice --- inherited from the LCF tradition of 1972 --- has forty years of accumulated libraries behind it, but it ties every corpus to a single $\LFnd$-point forever.

\Verum{}~\cite{VerumRepo,VerumDocs} breaks the tie. Its standard library is not a foundation; it is a \emph{classifier apparatus} that hosts foundations as external inputs. Framework axioms (\texttt{@framework(name, "citation")}) supply $\LFnd$-points explicitly; the \texttt{core.theory\_interop} module provides a generic Yoneda / Kan / descent machinery applicable to any theory; the \texttt{verum\_kernel} trusted base operates uniformly across foundations. A user importing \Verum{} does not import a foundation; she imports a mechanism for classifying foundations. This is the MSFS programme~\cite{Sereda2026MSFS} realised in code.

The paper establishes this claim formally (four correspondence theorems, \S\ref{sec:correspondences}), demonstrates it compositionally (theory-interop as generic $\LCls$ machinery, \S\ref{sec:theory-interop}), and positions it competitively (the foundation-neutrality gap no extant system can close, \S\ref{sec:competitive}).

\subsection{The refactor as architectural orthodoxy}
\label{subsec:refactor-orthodoxy}

An early \Verum{} contained a \texttt{core.mathesis} module specialised to the Mathesis meta-mathematical programme, and documentation positioned UHM (Unitary Holonomic Monism) as ``Verum's flagship proof corpus.'' Both commitments have been retired. The rename of \texttt{core.mathesis} to \texttt{core.theory\_interop} --- and the replacement of the UHM-specific documentation page with a generic \texttt{proof-corpora.md} --- are not cosmetic; they reflect an MSFS-structural principle that the paper articulates as the main thesis.

\begin{proposition}[Classifier neutrality under MSFS]
\label{prop:classifier-neutrality}
A member $\mathbf{F} \in \LCls$ is required by (M5) of~\cite[Definition \texttt{def:meta}]{Sereda2026MSFS} to be \emph{metatheory-parametrised}: $\mathbf{F}$ must resolve against an external Rich-metatheory parameter $S \in \RS$, not against one fixed internally. Consequently, any implementation aspiring to $\LCls$ membership must avoid privileging a specific $\LFnd$-point or application corpus.
\end{proposition}

\begin{proof}
(M5) is a defining condition of $\LCls$. Violating it means the candidate classifier fails (M5), hence is not in $\LCls$.
\end{proof}

The refactor enforces Proposition \ref{prop:classifier-neutrality} operationally: \Verum{}'s stdlib is now (M5)-compliant by construction. Specific applications (UHM, HoTT book, physical corpora, genomics ontologies) are external cogs that \emph{consume} \Verum{}, not sub-modules \Verum{} privileges.

\subsection{Contribution}
\label{subsec:contribution}

\begin{enumerate}[label=(C\arabic*)]
\item \textbf{Four correspondence theorems} (\S\ref{sec:correspondences}):
    \begin{itemize}[nosep]
    \item \texttt{@framework(...)} realises an indexing functor $\mathrm{Fw} : \mathrm{FwReg}(\Verum) \to \LFnd$ (Theorem \ref{thm:fnd-correspondence});
    \item Every \texttt{protocol} realises an $\LCls$-classifier with (M1)--(M5) data (Theorem \ref{thm:cls-correspondence});
    \item The trusted kernel realises T-2f* depth-stratification and inherits universal Yanofsky-immunity (Theorem \ref{thm:kernel-paradox-immune});
    \item The nine-strategy verification ladder refines the Diakrisis $\nu$-invariant (Theorem \ref{thm:verify-ladder}).
    \end{itemize}

\item \textbf{Domain-neutrality theorem} (\S\ref{sec:theory-interop}): the \texttt{core.theory\_interop} module realises generic Yoneda / Kan / descent operations uniformly across all \Verum{}-hosted foundations (Theorem \ref{thm:interop-domain-neutral}).

\item \textbf{Intensional self-positioning}: via the \Eff{}-topos invariant $\tau$ (MSFS \S 8.3), \Verum{} lies inside the $\MLTT$-slice of $\fM$, explaining the asymmetry of export compatibility (\S\ref{sec:mltt-ett}).

\item \textbf{Corpus-neutrality}: proof corpora are external \Verum{} cogs whose MSFS coordinate is computed, not assumed (\S\ref{sec:corpora}).

\item \textbf{Competitive analysis}: no extant proof assistant satisfies classifier-neutrality. The gap is architectural, not engineering (\S\ref{sec:competitive}).

\item \textbf{Knowledge-processing subsumption} (\S\ref{sec:kr}): the entire landscape of knowledge-representation paradigms --- \OWL{} description logics, \SWRL{} rules, modal / epistemic / temporal / dynamic / $\mu$-calculus, default / circumscription / answer-set, probabilistic / fuzzy / possibilistic --- is hosted in \Verum{} as external \texttt{@framework}-packages, each with its own MSFS coordinate. Industrial OWL2 reasoners (\Pellet, \HermiT, \FactPP, \ELK, \Konclude) are tableau engines without formal correctness; \Verum{} subsumes their decidable fragments with kernel-verified certification and extends reasoning across paradigms via \texttt{core.theory\_interop} gauge-transformations --- both unavailable elsewhere.
\end{enumerate}

\subsection{Competitive map}
\label{subsec:competitive-map}

Table~\ref{tab:competitive} summarises MSFS-level capabilities. The blank cell in \Verum{}'s row is the substance of this paper.

\begin{table}[ht]
\centering
\small
\caption{MSFS-level capabilities of major proof assistants. \checkmark = structurally realised; $\circ$ = realised via library / tactic; $\times$ = not present. Abbreviations: \emph{RT} = refinement types, \emph{DT} = dependent types, \emph{HoTT} = cubical HoTT / univalence, \emph{SMT} = first-class SMT backend, \emph{FA} = framework axioms with citation, \emph{K\textsubscript{5K}} = LCF-style kernel $<$5\,000 LOC target, \emph{9S} = nine-strategy verification dispatch, \emph{FN} = foundation-neutral (classifier rather than foundation).}
\label{tab:competitive}
\begin{tabular}{lcccccccc}
\toprule
System & RT & DT & HoTT & SMT & FA & K\textsubscript{5K} & 9S & FN \\ \midrule
\Verum{}       & \checkmark & \checkmark & \checkmark & \checkmark & \checkmark & \checkmark & \checkmark & \checkmark \\
\Coq/\textsc{Rocq} & $\times$   & \checkmark & $\circ$    & $\circ$    & $\times$   & $\circ$    & $\times$   & $\times$ \\
\Lean 4        & $\times$   & \checkmark & $\circ$    & $\circ$    & $\times$   & $\times$   & $\times$   & $\times$ \\
\Agda{}        & $\times$   & \checkmark & \checkmark & $\times$   & $\times$   & $\circ$    & $\times$   & $\times$ \\
\Isabelle/HOL  & $\times$   & $\circ$    & $\times$   & $\circ$    & $\times$   & \checkmark & $\times$   & $\times$ \\
\Fstar{}       & \checkmark & \checkmark & $\times$   & \checkmark & $\times$   & $\times$   & $\times$   & $\times$ \\
Dafny          & \checkmark & $\times$   & $\times$   & \checkmark & $\times$   & $\times$   & $\times$   & $\times$ \\
Liquid Haskell & \checkmark & $\times$   & $\times$   & \checkmark & $\times$   & $\times$   & $\times$   & $\times$ \\
Mizar          & $\times$   & $\circ$    & $\times$   & $\times$   & $\times$   & $\times$   & $\times$   & $\times$ \\
Metamath       & $\times$   & $\times$   & $\times$   & $\times$   & $\times$   & \checkmark & $\times$   & $\times$ \\
\bottomrule
\end{tabular}
\end{table}

The FN column is the new post-refactor invariant --- the one no competitor has because no competitor architecturally separates language from foundation. The other seven columns follow from public documentation of each system.

\subsection{Paper organisation}
\label{subsec:organisation}

\S\ref{sec:msfs-recalled} recalls MSFS strata. \S\ref{sec:verum-architecture} summarises \Verum{}'s architecture after the classifier-neutrality refactor. \S\ref{sec:correspondences} proves the four correspondence theorems. \S\ref{sec:theory-interop} establishes domain-neutrality of theory interop. \S\ref{sec:mltt-ett} fixes \Verum{}'s intensional coordinate. \S\ref{sec:corpora} discusses proof corpora as external $\LFnd$-consumers. \S\ref{sec:kr} develops knowledge-processing subsumption: \OWL{} and the wider KR landscape as \texttt{@framework}-hosted foundations, with certification dominance over industrial reasoners. \S\ref{sec:competitive} delivers the strict qualitative and structural competitive analysis across proof assistants, OWL2 reasoners, and HOL-based ontologies. \S\ref{sec:roadmap} sketches a dominance roadmap. \S\ref{sec:conclusion} concludes.

\section{MSFS stratification recalled}
\label{sec:msfs-recalled}

We recall from~\cite{Sereda2026MSFS} the four strata and their defining conditions, sufficient for self-contained reading.

\begin{definition}[Rich-metatheory]
\label{def:rs-recalled}
A metatheory $S$ is a \emph{Rich-foundation}, written $S \in \RS$, if it satisfies (R1) Peano interpretation, (R2) r.e.-axiomatisation, (R3) internal metatheory of sufficient strength, (R4) totality of recursive functions, and (R5) Morita-stability of the admissible model class.
\end{definition}

\begin{definition}[The four strata]
\label{def:strata-recalled}
\begin{align*}
\LFnd     &= \{ S : S \in \RS \} \quad \text{(Rich-foundations; e.g.\ } \ZFC, \HoTT, \CIC, \text{NCG}, \Eff, \text{etc.)}. \\
\LCls     &= \{ \mathbf{F} : \mathbf{F} \text{ satisfies (M1)--(M5)} \} \quad \text{(classifiers)}. \\
\LClsMax  &= \{ \mathbf{F} \in \LCls : \mathbf{F} \text{ also satisfies (Max-1)--(Max-4)} \} \quad \text{(maximal classifiers)}. \\
\LAbs     &= \{ X : X \text{ satisfies } (F_S) \wedge (\Pi_{4,S,n}) \wedge (\Pi_{3\text{-max},S,n}) \} \quad \text{(absolute foundations)}.
\end{align*}
AFN-T (MSFS Theorem 5.1) establishes $\LAbs = \emptyset$, with uniform five-axis absoluteness in $(S, n, \mu, \xi, \pi)$.
\end{definition}

The conditions defining $\LCls$ are:
\begin{description}[nosep]
\item[(M1)] Locally small 2-category.
\item[(M2)] Classification functor $\mathrm{Cl}_{\mathbf{F}} : \mathbf{F} \to \fM$.
\item[(M3)] Accessible reflection.
\item[(M4)] Non-generative (organises, does not derive).
\item[(M5)] Metatheory-parametrised.
\end{description}

\begin{convention}[Diakrisis extensions]
\label{conv:diakrisis-extensions}
Three constructions from the Diakrisis programme~\cite{Diakrisis} are invoked:
\begin{enumerate}[label=(D\arabic*), nosep]
\item $\nu$-invariant: each articulation $\alpha$ carries an ordinal rank $\nu(\alpha)$ measuring its $\mathsf{M}$-iteration depth (23.T1).
\item T-2f*: depth-stratification axiom blocking the Yanofsky family of self-referential paradoxes.
\item Maximality proofs 103.T--106.T: establish $\LClsMax \neq \emptyset$ with Diakrisis as theorem-level witness.
\end{enumerate}
\end{convention}

\section{\Verum{} architecture: classifier-neutrality by construction}
\label{sec:verum-architecture}

\Verum{} is a statically-typed compiled systems language with three reserved keywords (\texttt{let}, \texttt{fn}, \texttt{is}); all further syntax is contextual. The architecture after the classifier-neutrality refactor has six layers, each domain-neutral.

\subsection{Core language}
\label{subsec:core-language}

\begin{enumerate}[nosep]
\item \textbf{Refinement types} $\tau\{\varphi(\mathtt{self})\}$: SMT-discharged subtypes.
\item \textbf{Dependent types}: $\Sigma$ and $\Pi$ with cubical path equality.
\item \textbf{Cubical path types} with CCHM normalisation~\cite{CCHM2015}: \texttt{Path<A>(a,b)}, \texttt{hcomp}, \texttt{transp}, \texttt{Glue}, computational univalence.
\item \textbf{Protocols with associated types}: a protocol is a $\Pi$-indexed record of methods and types with specialisation.
\end{enumerate}

None of these features commits to any Rich-foundation; all are encodable in any R-S that satisfies (R1)--(R5).

\subsection{The trusted kernel}
\label{subsec:kernel}

\texttt{verum\_kernel} defines an explicit \texttt{CoreTerm} calculus with $\Sigma$, $\Pi$, \texttt{Path}, \texttt{hcomp}, \texttt{transp}, \texttt{Glue}, inductive types, universe levels, refinement subtypes, and framework-axiom constructors. Every tactic, every SMT certificate, every elaborator output is coerced to a \texttt{CoreTerm} and re-checked. Target size on completion: $<$5\,000 LOC. Current: 1\,180 LOC.

The kernel does not hard-code any foundation. Universe levels are indexed by natural numbers; framework axioms are first-class \texttt{CoreTerm::Axiom(name, type)} values with no privileged subset.

\subsection{Framework axioms (\texttt{core.math.frameworks})}
\label{subsec:framework-axioms}

The stdlib ships six example framework packages: Lurie HTT, Schreiber DCCT, Connes reconstruction, Petz classification, Arnold--Mather, Baez--Dolan (together 36 citation-tagged axioms). These are \emph{examples}, not privileged: users add their own via \texttt{@framework(name, "citation")} on any \texttt{axiom} declaration.

\begin{lstlisting}[caption={User-defined framework axiom.}]
@framework(my_programme, "Smith 2026 Theorem 3.2")
axiom my_external_result<T>(x: T) -> Bool
    ensures x.my_predicate();
\end{lstlisting}

The command \verb|verum audit --framework-axioms| enumerates the full dependency list of any proof corpus. This datum is the corpus's MSFS-coordinate (see \S\ref{sec:corpora}).

\subsection{The nine-strategy verification ladder}
\label{subsec:verification-ladder}

Each function carries a \texttt{@verify(strategy)} attribute selecting one of nine strategies:
\[
\texttt{runtime, static, fast, formal, proof, thorough, reliable, certified, synthesize}.
\]
These are \emph{intents}, not backends. The \texttt{verum\_smt} router picks an SMT backend (Z3, CVC5, E, Vampire) by goal shape under the chosen intent. The kernel records the strategy in the proof term, guaranteeing replay at the demanded strength.

\subsection{Theory interoperation (\texttt{core.theory\_interop})}
\label{subsec:theory-interop-intro}

The \texttt{core.theory\_interop} stdlib module (formerly \texttt{core.mathesis}) implements three fundamental operations in the computational $\infty$-topos of formally-represented theories:
\begin{enumerate}[nosep]
\item \texttt{load\_theory(T)} --- Yoneda embedding $T \mapsto y(T)$;
\item \texttt{translate(T\textsubscript{1}, T\textsubscript{2}, partial)} --- Kan extension of a partial functor;
\item \texttt{check\_coherence(translations)} --- descent verification on the \v{C}ech nerve.
\end{enumerate}
The module is not specialised to any programme. Any theory expressible against the \texttt{Theory} and \texttt{Claim} interfaces is a valid input. \S\ref{sec:theory-interop} establishes this domain-neutrality formally.

\subsection{User-facing verification API (\texttt{core.verify}, \texttt{core.proof})}
\label{subsec:user-verify}

Two stdlib layers expose verification to user code:
\begin{itemize}[nosep]
\item \texttt{core.verify} --- level / attempt / certificate types; \texttt{Verified<T>}, \texttt{Proven<P>}, level-branching.
\item \texttt{core.proof} --- proof-carrying code bundles (\texttt{pcc.vr}, Necula-style), reflection, and the tactic library.
\end{itemize}
Both are domain-neutral: they operate against whatever foundation the user's \texttt{@framework} imports have supplied.

\section{Four correspondence theorems}
\label{sec:correspondences}

We establish the four correspondences. Each correspondence is stated in a form free of specific-instance bias: the theorems quantify over all \Verum{} users, not over a privileged application.

\subsection{Correspondence I: framework axioms $\leftrightarrow$ $\LFnd$}
\label{subsec:corr1}

\begin{theorem}[Framework-axiom registry as $\LFnd$-indexing functor]
\label{thm:fnd-correspondence}
Let $\mathrm{FwReg}(\Verum)$ denote the \texttt{AxiomRegistry} of a \Verum{} installation (stdlib + user cogs), with morphisms given by \texttt{@framework}-citation inclusions. There exists a faithful functor
\[
\mathrm{Fw} : \mathrm{FwReg}(\Verum) \longrightarrow \LFnd
\]
sending each \texttt{@framework(name, "cit")} declaration to the Rich-foundation of the cited theorem. The six stdlib-shipped packages embed as:
\begin{align*}
\texttt{lurie\_htt}           &\;\longmapsto\; \text{Lurie } (\infty, 1)\text{-topos theory} \in \LFnd, \\
\texttt{schreiber\_dcct}      &\;\longmapsto\; \text{cohesive } (\infty, 1)\text{-topos theory} \in \LFnd, \\
\texttt{connes\_reconstruction} &\;\longmapsto\; \text{noncommutative geometry} \in \LFnd, \\
\texttt{petz\_classification} &\;\longmapsto\; \text{operator-algebraic quantum information} \in \LFnd, \\
\texttt{arnold\_mather}       &\;\longmapsto\; \text{smooth singularity theory} \in \LFnd, \\
\texttt{baez\_dolan}          &\;\longmapsto\; (\infty, n)\text{-categorical coherence} \in \LFnd.
\end{align*}
These are \emph{examples}, not privileged: $\mathrm{Fw}$ is defined uniformly on user-supplied \texttt{@framework} declarations without distinguishing stdlib from non-stdlib entries.
\end{theorem}

\begin{proof}
Each \texttt{@framework(name, "cit")} declaration asserts a (R1)--(R5)-satisfying metatheoretic assumption. Faithfulness: distinct \texttt{@framework} identifiers with different citations induce distinct $\LFnd$-points because the underlying metatheories differ in at least one (R1)--(R5) datum. Functoriality: citation-inclusion morphisms (cases where one framework is a sub-collection of another) are preserved. The functor makes no reference to stdlib versus user-cog provenance of the declaration.
\end{proof}

\begin{corollary}[Audit command as $\LFnd$-classifier]
\label{cor:audit-classifier}
\verb|verum audit --framework-axioms| on any proof corpus $P$ returns the image $\mathrm{Fw}(P) \subseteq \LFnd$ --- the corpus's MSFS-coordinate as a computable set of $\LFnd$-points.
\end{corollary}

\subsection{Correspondence II: protocols $\leftrightarrow$ $\LCls$}
\label{subsec:corr2}

\begin{theorem}[Protocols as $\LCls$-classifiers]
\label{thm:cls-correspondence}
Let $\mathbf{P}$ be a \Verum{} \texttt{protocol} declaration and $\cF_{\mathbf{P}}$ the 2-category of types implementing $\mathbf{P}$ (objects: types \texttt{T} with \texttt{implement P for T}; 1-morphisms: protocol-preserving functions; 2-morphisms: path equalities). Then $\cF_{\mathbf{P}}$ satisfies (M1)--(M5):
\begin{description}[nosep]
\item[(M1)] Locally small, indexed by type variables.
\item[(M2)] Classification functor $\Cl_{\cF_{\mathbf{P}}} : \cF_{\mathbf{P}} \to \fM$ via gauge quotient.
\item[(M3)] Accessible via \texttt{@derive}-based canonical method synthesis.
\item[(M4)] Non-generative: protocols receive axioms from implementations; they do not derive them.
\item[(M5)] Metatheory-parametrised: protocols resolve against the ambient $S \in \RS$ imported by \texttt{@framework} declarations in scope.
\end{description}
(M5) is precisely classifier-neutrality: the protocol does not fix its Rich-foundation; it accepts whatever the enclosing scope supplies.
\end{theorem}

\begin{proof}
(M1) by construction of \Verum{}'s type-class-resolution index. (M2) by gauge-quotient of $\cF$ factored through protocol-preserving morphisms. (M3) accessibility follows from finite-arity restrictions on protocol method signatures. (M4) by the grammatical separation of \texttt{protocol} (interface) and \texttt{implement} (witness supply). (M5) by the module system's mount-scope resolution of \texttt{@framework} imports: the ambient R-S is an external lookup, not a protocol-internal commitment.
\end{proof}

\begin{example}
Stdlib protocols such as \texttt{Functor}, \texttt{Monad}, \texttt{Category}, \texttt{Topos} are classifiers in $\LCls$. The \texttt{core.math.cohesive} module is a full \Verum{}-internal realisation of Schreiber's cohesive framework~\cite{Schreiber2013} --- one of the three partial classifiers established by MSFS Theorem \texttt{thm:meta-mult}.
\end{example}

\begin{remark}[No single protocol reaches $\LClsMax$]
\label{rem:no-max-at-protocol-level}
No individual \Verum{} protocol satisfies all of (Max-1)--(Max-4): no single protocol classifies every Rich-foundation. Reaching $\LClsMax$ requires the composite of the protocol-family together with the kernel's depth-stratification and the framework-axiom registry (see \S\ref{sec:roadmap}, R1).
\end{remark}

\subsection{Correspondence III: kernel $\leftrightarrow$ T-2f* + universal paradox-immunity}
\label{subsec:corr3}

\begin{theorem}[Kernel realises T-2f* and inherits Yanofsky-immunity]
\label{thm:kernel-paradox-immune}
Let $\mathrm{Kernel}(\Verum)$ be the closed subcategory of \texttt{CoreTerm}s accepted by \texttt{verum\_kernel::infer} / \texttt{verify\_full}. Then:
\begin{enumerate}[label=(\alph*), nosep]
\item Every \texttt{CoreTerm} carries a computed $\mathsf{M}$-depth $d \in \mathrm{Ord}$ equal to its universe-level plus type-constructor nesting depth.
\item Comprehension in the kernel (\texttt{CoreTerm::Set}, \texttt{CoreTerm::Refinement}) succeeds only when the comprehension predicate has $\mathsf{M}$-depth strictly less than the constructed object's depth --- the T-2f* axiom realised operationally.
\item Every self-referential term that reduces to a Yanofsky universal diagonal~\cite{Yanofsky2003} is rejected by the kernel. \Verum{} inherits universal paradox-immunity from Diakrisis Theorem 105.T.
\end{enumerate}
\end{theorem}

\begin{proof}
(a) \texttt{CoreTerm::Universe} encodes level explicitly; $\Sigma$, $\Pi$, \texttt{Eq}, \texttt{Path} constructors increment it. Nesting depth is a syntactic recursion over the term.

(b) Rule \texttt{K-Refine} of the kernel's typing judgement checks the predicate depth. Any comprehension whose predicate exceeds the object level fails elaboration and the kernel rejects the resulting term.

(c) By Diakrisis 105.T, every Yanofsky-reducible paradox reduces to a weakly point-surjective $\alpha : Y \to T^Y$ in cartesian-closed context with $\mathrm{dp}(T^Y) = \mathrm{dp}(Y)+1$. The depth-strict condition in (b) blocks exactly this: no kernel-accepted $\alpha$ can witness weak point surjectivity while preserving depth. Transporting the Diakrisis proof of 105.T to the kernel's typing judgement yields (c).
\end{proof}

\begin{corollary}[\Verum{} is universally paradox-immune]
\label{cor:verum-universal-immunity}
No proof corpus accepted by the kernel contains a Yanofsky-reducible paradoxical term. The immunity holds for every such paradox, not only the five historical families (Russell, Curry, Grelling, Burali-Forti, Girard).
\end{corollary}

\begin{remark}[Contrast with Coq / Lean universe hierarchies]
\label{rem:universe-comparison}
\Coq{} and \Lean{} implement universe levels with cumulativity rules that partially realise depth-stratification. The decisive difference: in \Coq{} and \Lean{} universe consistency is an \emph{engineering outcome}; in \Verum{} it is a \emph{theorem} about the kernel via Diakrisis 105.T. The transport from engineering claim to theorem is exactly what MSFS-native design buys.
\end{remark}

\subsection{Correspondence IV: verification strategies $\leftrightarrow$ $\nu$-invariant}
\label{subsec:corr4}

\begin{theorem}[Verification ladder refines $\nu$]
\label{thm:verify-ladder}
The nine verification strategies monotonically refine the Diakrisis $\nu$-invariant on proof-goal depth:
\begin{center}
\begin{tabular}{ll}
\toprule
Strategy & $\nu$-image \\ \midrule
\texttt{runtime} & $\nu = 0$ (execution-time assertion) \\
\texttt{static}  & $\nu < \omega$ (finite dataflow) \\
\texttt{fast}    & $\nu < \omega$ (finite SMT, possibly incomplete) \\
\texttt{formal}  & $\nu = \omega$ (full SMT discharge) \\
\texttt{proof}   & $\nu = \omega$ (alias) \\
\texttt{thorough}  & $\nu = \omega \cdot 2$ (parallel multi-strategy) \\
\texttt{reliable}  & $\nu = \omega \cdot 2$ (alias) \\
\texttt{certified} & $\nu = \omega \cdot 2$ + cross-kernel export witness \\
\texttt{synthesize} & $\nu \leq \omega \cdot 3 + 1$ (inverse search in $\fM$) \\
\bottomrule
\end{tabular}
\end{center}
Higher-strategy claims inhabit the closure of lower-strategy claims.
\end{theorem}

\begin{proof}
Finite-step strategies (\texttt{runtime} / \texttt{static} / \texttt{fast}) correspond to $\nu < \omega$. \texttt{formal} discharges arbitrarily-bounded-arity SMT goals, $\nu = \omega$. \texttt{thorough} and its alias \texttt{reliable} close under method choice via multi-strategy racing, advancing one $\mathsf{M}$-layer, $\nu = \omega \cdot 2$. \texttt{certified} adds a cross-kernel witness at the same $\nu$. \texttt{synthesize} performs inverse-search across the gauge-quotient $\fM$, reaching the maximum $\nu$ attainable by any corpus \Verum{} hosts.
\end{proof}

\section{Domain-neutrality of theory interoperation}
\label{sec:theory-interop}

The \texttt{core.theory\_interop} module is \Verum{}'s most technically ambitious stdlib component. The refactor renamed it from \texttt{core.mathesis} to signal that it is not tied to the Mathesis programme. The theorem below formalises this: the module's three operations are uniform across all valid theory inputs, regardless of domain.

\begin{theorem}[Domain-neutrality of theory interoperation]
\label{thm:interop-domain-neutral}
Let $\mathcal{T}$ be the class of theories expressible against the \texttt{core.math.epistemic::Theory} interface. The three operations
\begin{align*}
\texttt{load\_theory} &: \mathcal{T} \to \mathrm{LoadedTheory}, \\
\texttt{translate}    &: \mathrm{LoadedTheory} \times \mathrm{LoadedTheory} \times \mathrm{Partial} \to \mathrm{Translation}, \\
\texttt{check\_coherence} &: \mathrm{Registry} \times \mathrm{Pairs} \to \mathrm{CoherenceResult}
\end{align*}
are functors that commute with every automorphism $\sigma : \mathcal{T} \to \mathcal{T}$ preserving the \texttt{Theory} interface. In particular, none of the three privileges any specific subclass of $\mathcal{T}$ (physical theories, mathematical theories, engineering ontologies, etc.).
\end{theorem}

\begin{proof}
The operations are defined entirely in terms of the \texttt{Theory} / \texttt{Claim} / \texttt{EpistemicTopology} interfaces (\texttt{core.math.epistemic}). Any automorphism of $\mathcal{T}$ preserving these interfaces commutes with the operations because they are parametric in the interface data.
\end{proof}

\begin{corollary}[Theory interop as generic classifier]
\label{cor:interop-generic-classifier}
\texttt{core.theory\_interop} together with the kernel realises a $\LCls$ classifier that accepts any theory in $\mathcal{T}$ as input. The classifier is (M5)-compliant by Theorem \ref{thm:interop-domain-neutral}: the ambient R-S is the input theory's own metatheory, not an internally fixed one.
\end{corollary}

\begin{remark}[The refactor payoff]
\label{rem:refactor-payoff}
Pre-refactor, the module was tied to Mathesis theories specifically, which would have failed (M5) and violated classifier-neutrality (Proposition \ref{prop:classifier-neutrality}). The rename and decoupling restore (M5)-compliance and make the correspondences I--IV hold without reference to any specific application.
\end{remark}

\section{\Verum{}'s intensional coordinate: the $\MLTT$-slice}
\label{sec:mltt-ett}

\Verum{} implements cubical path types natively (\texttt{core.math.hott}, \texttt{core.math.cubical}) with CCHM normalisation and propositional identity. It does not implement definitional-equality reflection, which would place it in the $\ETT$ slice~\cite{Hofmann1995}.

\begin{proposition}[\Verum{} lies in the $\MLTT$-slice]
\label{prop:verum-in-mltt-slice}
With respect to the intensional refinement functor $\II : \cF^{\mathrm{op}} \to \Sint$ of MSFS Theorem 8.1:
\[
\tau(\II(\Verum)) = 1,
\]
where $\tau$ is the \Eff{}-internal computability invariant. Equivalently: \Verum{}'s type-check problem admits an \Eff{}-internal normaliser.
\end{proposition}

\begin{proof}
\Verum{}'s type-check is the strong-normalisation checker of CCHM cubical type theory extended with SMT-side goals. CCHM is strongly normalising~\cite{Huber2019}; SMT certificates are themselves \texttt{CoreTerm}s rechecked in the kernel. The kernel's acceptance predicate is total in $\Eff$, hence $\tau(\II(\Verum)) = 1$.
\end{proof}

\begin{corollary}[Export compatibility asymmetry]
\label{cor:export-compat}
Certified export (\verb|verum export --to|) is structurally lossless to \Coq, \Lean, \Agda, Dedukti, Metamath --- all $\MLTT$-slice systems. Export to $\ETT$-slice systems (\Isabelle/HOL) crosses an intensional fibre of $\fM$: $\tau$ drops from $1$ to $0$, witnessing foundational asymmetry. This structurally explains the folklore observation that \Coq{} $\to$ \Isabelle{} export is harder than the reverse.
\end{corollary}

\section{Proof corpora as external $\LFnd$-consumers}
\label{sec:corpora}

A \Verum{} proof corpus is an external cog (package) satisfying a generic pattern documented in~\cite{VerumCorpora}. It is \emph{not} a stdlib extension. This is the architectural counterpart of Proposition \ref{prop:classifier-neutrality}: corpora consume \Verum{}, they do not privilege it.

\begin{definition}[Proof corpus pattern]
\label{def:corpus-pattern}
A \emph{\Verum{} proof corpus} is a cog providing:
\begin{enumerate}[nosep]
\item A \texttt{verum.toml} research-profile manifest declaring default \texttt{@verify} strategy, timeouts, and allowed framework-axiom set.
\item A layered \texttt{src/} tree: foundations $\to$ combinatorial $\to$ analytical $\to$ categorical $\to$ core-theorems $\to$ fundamental-closures.
\item A \texttt{certificates/} sub-tree populated by \verb|verum export --to {coq, lean, dedukti, metamath}|.
\item An auto-regenerated \texttt{theorem\_index.md} status map (theorem $\to$ stratum $\to$ \texttt{@verify}-strategy $\to$ certificate checksum).
\end{enumerate}
\end{definition}

\begin{theorem}[Corpus MSFS coordinate]
\label{thm:corpus-msfs-coordinate}
Any \Verum{} proof corpus $P$ conforming to Definition \ref{def:corpus-pattern} has a computable MSFS coordinate
\[
\mathrm{Coord}(P) = \bigl(\, \mathrm{Fw}(P),\; \{ \nu(\theta) : \theta \in P \},\; \tau(\II(P)) \,\bigr),
\]
where $\mathrm{Fw}$ is the framework functor (Theorem \ref{thm:fnd-correspondence}), $\nu(\theta)$ the Diakrisis invariant (Theorem \ref{thm:verify-ladder}), and $\tau$ the intensional invariant (Proposition \ref{prop:verum-in-mltt-slice}).
\end{theorem}

\begin{proof}
$\mathrm{Fw}(P)$ is computed by \verb|verum audit --framework-axioms|. The per-theorem $\nu(\theta)$ is read from the \texttt{@verify} attribute via Theorem \ref{thm:verify-ladder}. $\tau(\II(P))$ is inherited from \Verum{}'s slice position (Proposition \ref{prop:verum-in-mltt-slice}) because all corpora run through the same kernel.
\end{proof}

\begin{corollary}[Corpus neutrality]
\label{cor:corpus-neutrality}
No specific corpus is privileged by \Verum{}. Any cog meeting Definition \ref{def:corpus-pattern} receives an MSFS coordinate on the same footing as any other. The architecture does not distinguish mathematical, physical, engineering, or cross-domain corpora.
\end{corollary}

This is the operational meaning of the refactor. The former commitment to UHM as ``\Verum{}'s flagship corpus'' is incompatible with Corollary \ref{cor:corpus-neutrality}: a host that flags one corpus violates classifier-neutrality. The refactor eliminated the privilege; the resulting architecture is the one the paper describes.

\section{Knowledge Processing: Ontologies and Reasoning as MSFS-Hosted Foundations}
\label{sec:kr}

The Semantic Web standard \OWL~\cite{OWL2Primer,OWL2Spec} defines a description-logic fragment (\textsf{SROIQ}) whose industrial reasoners --- \Pellet~\cite{Pellet2007}, \HermiT~\cite{HermiT2014}, \FactPP~\cite{FactPP2006}, \ELK~\cite{ELK2014}, \Konclude~\cite{Konclude2014} --- provide tableau-based consistency, classification, and instance-checking. Classical upper-ontology projects --- \SUMO~\cite{SUMO2004}, \Cyc~\cite{Cyc1995}, \DOLCE~\cite{DOLCE2003}, \BFO~\cite{BFO2015} --- embed knowledge in higher-order classical logic. Between them lies a structured zoo: \SWRL{} rules~\cite{SWRL2004}, Common Logic (ISO 24707), \RDF / \textsc{RDF\,S}, modal / epistemic / temporal / dynamic logics, default / circumscription / answer-set programming, probabilistic / fuzzy / possibilistic extensions.

In MSFS terms, each paradigm is a point of $\LFnd$. The totality forms a sub-stack $\fM^{\mathrm{KR}} \subset \fM$; each entry satisfies (R1)--(R5) of Definition~\ref{def:rs-recalled} with a specific decidability-by-fragment profile and categorical semantics. This section establishes that \Verum{} is the first system to host this entire sub-stack uniformly, with kernel-verified certification on decidable fragments and interactive extension to undecidable ones --- a capability structurally unavailable to every industrial reasoner and every existing proof assistant.

\subsection{The knowledge-representation landscape as $\LFnd$ sub-stack}
\label{subsec:kr-landscape}

\begin{definition}[The KR sub-stack $\fM^{\mathrm{KR}}$]
\label{def:kr-substack}
Let $\fM^{\mathrm{KR}} \subset \fM$ be the full sub-$2$-stack of Rich-foundations specialising to knowledge representation. Exhibited points include:
\begin{itemize}[nosep]
\item \textbf{Description logics}: \textsf{AL}, \textsf{ALC}, \textsf{SHIQ}, \textsf{SHOIQ}, \textsf{SROIQ} ($=$ \OWL-\textsc{DL}), \textsf{EL} ($=$ \OWL-\textsc{EL}), \textsf{DL-Lite} ($=$ \OWL-\textsc{QL}), Datalog fragment ($=$ \OWL-\textsc{RL}).
\item \textbf{Rule languages}: \SWRL, RIF, Datalog, stable-model / answer-set programs.
\item \textbf{Higher-order ontologies}: HOL shallow embeddings (\SUMO, \Cyc, \DOLCE, \BFO), second-order logic with standard semantics.
\item \textbf{Modal families}: K, T, S4, S5, KD45, GL (provability), propositional dynamic logic (PDL), modal $\mu$-calculus.
\item \textbf{Epistemic / doxastic logics}: Hintikka's S5-based epistemic systems, multi-agent $\mathrm{K}_i$-calculi.
\item \textbf{Temporal logics}: LTL, CTL, CTL$^{\ast}$, ATL, metric-temporal MTL.
\item \textbf{Non-monotonic reasoning}: default logic (Reiter), circumscription (McCarthy), autoepistemic logic (Moore), answer-set programming.
\item \textbf{Uncertainty-bearing}: Bayesian networks, P-\textsf{SROIQ} / P-\textsf{SHOQ} probabilistic DL, fuzzy DL (Łukasiewicz, Gödel, product t-norms), possibilistic logic (Dubois--Prade).
\item \textbf{Common Logic} (ISO 24707), \textbf{Cyc\,L} variable-arity predicates.
\end{itemize}
\end{definition}

\begin{proposition}[KR-landscape is an $\LFnd$ sub-stack]
\label{prop:kr-landscape}
Each paradigm in Definition~\ref{def:kr-substack} satisfies (R1)--(R5) of Definition~\ref{def:rs-recalled}, hence inhabits $\LFnd$. The totality $\fM^{\mathrm{KR}}$ is Morita-closed under standard translations (DL-to-FOL standard translation, HOL-to-DT Curry--Howard, modal-to-FOL Blackburn--de\,Rijke--Venema).
\end{proposition}

\begin{proof}[Sketch]
Each paradigm provides: (R1) Peano interpretation via its arithmetic sub-fragment (trivially for FOL-reducible DLs; through HOL for upper ontologies); (R2) r.e.-axiomatisation (finite or decidably r.e.\ axiom schemata in every listed case); (R3) standard model-class in a metatheory $S$ (ZFC for classical DL, Kripke-frame models for modal families, probability space models for probabilistic extensions); (R4) recursive totality of well-formed-formula checkers; (R5) Morita-stability via the standard translation zoo. Morita-closure of $\fM^{\mathrm{KR}}$ under the cited translations is well-known in description-logic theory~\cite{BaaderHandbook2017}.
\end{proof}

\subsection{OWL\,2 as \texttt{@framework}-hosted foundation}
\label{subsec:owl2-hosted}

The cited HOL formalisation of \OWL-FS operators~\cite{RudolphHorrocksHOL2019} shows every functional-syntax operator of \OWL{} --- \texttt{SubClassOf}, \texttt{ObjectIntersectionOf}, \texttt{ObjectSomeValuesFrom}, \texttt{DataHasValue}, \texttt{EquivalentClasses}, \texttt{TransitiveObjectProperty}, the full SROIQ operator set --- admits a definition in a simplest HOL language. Since \Verum's DT kernel embeds HOL shallowly (HOL $\hookrightarrow$ DT via refinement types plus proof-irrelevance discipline), each \OWL-FS operator descends to a native \Verum{} \texttt{CoreTerm}.

\begin{construction}[The \texttt{core.kr.owl2} stdlib package]
\label{constr:owl2-pkg}
Let \texttt{core.kr.owl2} be a stdlib cog containing \texttt{@framework(owl2\_dl, "W3C OWL\,2 Recommendation 2012")} axioms for the \textsf{SROIQ} signature, Kripke-style model-theoretic semantics, the standard translation to FOL, and certified tableau-algorithm implementation. Reasoning tasks are exposed as \texttt{@verify}-attributed functions whose soundness is a kernel-verified theorem relative to the W3C model-theoretic semantics.
\end{construction}

\begin{lstlisting}[caption={Illustrative skeleton of \texttt{core.kr.owl2}.}]
module core.kr.owl2

@framework(owl2_dl, "W3C OWL 2 Rec 2012, §9 SROIQ")
axiom Concept : Type;
axiom Role    : Type;
axiom Ind     : Type;  // individuals

@framework(owl2_dl, "W3C OWL 2 Rec 2012, §3 Class Expressions")
axiom Intersection : Concept -> Concept -> Concept;
axiom Union        : Concept -> Concept -> Concept;
axiom Complement   : Concept -> Concept;
axiom Exists       : Role -> Concept -> Concept;
axiom Forall       : Role -> Concept -> Concept;
axiom AtLeast      : Int -> Role -> Concept -> Concept;
axiom AtMost       : Int -> Role -> Concept -> Concept;
// ... full SROIQ operator set (nominals, inverse roles, transitive roles, etc.)

@framework(owl2_dl, "W3C OWL 2 Rec 2012, §9 Semantics")
axiom models : Interpretation -> Axiom -> Bool;

@verify(certified)
fn classify(ont: Ontology) -> ClassificationHierarchy
    ensures sound_wrt(owl2_dl, ont, result);

@verify(formal)
fn consistent(ont: Ontology) -> Bool
    ensures result = true <-> exists(I: Interpretation). models_all(I, ont);
\end{lstlisting}

\subsection{Certification dominance: strict and honest}
\label{subsec:kr-certification-dominance}

\begin{theorem}[Certified-reasoning dominance on the decidable KR fragment]
\label{thm:kr-cert-dominance}
Let $P$ be any reasoning task (consistency, classification, realisation, conjunctive query answering) expressible in \textsf{SROIQ} or any Morita-reducible DL fragment. Then:
\begin{enumerate}[label=(\alph*), nosep]
\item \Verum{} supplies an implementation $f_P : \mathrm{Input}(P) \to \mathrm{Output}(P)$ whose soundness against the W3C model-theoretic semantics is a kernel-verified theorem.
\item \Pellet, \HermiT, \FactPP, \ELK, \Konclude{} supply implementations whose soundness is established by empirical testing against the OWL\,2 W3C conformance suite, not by formal proof relative to their own source code.
\end{enumerate}
\end{theorem}

\begin{proof}
(a) By Construction~\ref{constr:owl2-pkg} each $P$ admits a \Verum{} implementation whose proof of soundness against \texttt{core.kr.owl2}'s \texttt{@framework(owl2\_dl, ...)} axioms is a \texttt{@verify(certified)}-attributed \texttt{CoreTerm}, kernel-checked by Theorem~\ref{thm:kernel-paradox-immune}. (b) Publicly documented testing methodologies of the named reasoners rely on the W3C conformance corpus~\cite{OWL2Tests} and internal regression testing; no formal soundness proof of the implementation is available for any of them.
\end{proof}

\begin{remark}[Speed parity as an engineering deliverable]
\label{rem:kr-cert-tradeoff}
Theorem~\ref{thm:kr-cert-dominance} is a structural claim about \emph{correctness guarantee}; speed is a separate engineering concern with no architectural ceiling on \Verum's side. The tableau algorithm is the same across implementations; \Pellet{} and \HermiT{} run it in Java, \FactPP{} and \Konclude{} in C++. \Verum{} owns its compiler, runtime, and kernel end-to-end (current codebase: $\sim 3{,}937$ LOC kernel, $\sim 107{,}000$ LOC SMT infrastructure across Z3 / CVC5 / refinement / exhaustiveness backends via a shipped \texttt{backend\_trait} / \texttt{backend\_switcher} plugin interface, $\sim 59{,}859$ LOC verification and tactics, and certificate emission to Dedukti / Coq / Lean / OpenTheory / Metamath). A dedicated Rust-native \textsf{SROIQ} tableau reasoner is therefore an additional backend within the existing plugin architecture: native tableau speed (at or below the per-step cost of current industrial reasoners, given Rust's performance profile), emitting \texttt{CoreTerm}-certificates that the kernel re-checks in a single polynomial-time final pass. \emph{Speed parity with} \Pellet{} / \HermiT{} \emph{is an engineering deliverable of R6; kernel-verified correctness is already architectural}. The two advantages compose rather than trade: a user chooses between \emph{an empirically-validated fast opinion} (\Pellet) and \emph{an equally-fast formally-verified certificate} (\Verum{} post-R6). On speed today, the generic-SMT \texttt{@verify(formal)} path is slower than specialised tableau; this reflects the current absence of \texttt{core.kr.owl2\_tableau}, not an architectural constraint on \Verum.
\end{remark}

\subsection{Cross-paradigm translation via \texttt{core.theory\_interop}}
\label{subsec:kr-cross-paradigm}

\begin{theorem}[Cross-paradigm KR translation with certified gauge-witnesses]
\label{thm:kr-cross-paradigm}
Let $K_1, K_2 \in \fM^{\mathrm{KR}}$ be two paradigms in Definition~\ref{def:kr-substack} and $\sigma : K_1 \to K_2$ a standard translation (DL-to-FOL, FOL-to-HOL, HOL-to-DT, modal-to-FOL, etc.). The \texttt{core.theory\_interop} operations instantiate $\sigma$ as a functorial translation whose Kan-extension coherence (Theorem~\ref{thm:interop-domain-neutral}) is kernel-verified in \Verum. Every ontology $O \in K_1$ thus produces an ontology $\sigma(O) \in K_2$ together with a gauge-transformation certificate in $\fM^{\mathrm{KR}}$.
\end{theorem}

\begin{proof}[Sketch]
Each standard translation is a functor of signatures; the \texttt{core.theory\_interop} module implements Yoneda / Kan / descent operations parametric over the \texttt{Theory} / \texttt{Claim} interfaces (Theorem~\ref{thm:interop-domain-neutral}), hence supports $\sigma$ uniformly. Kernel-verification of coherence reduces to coherence of the underlying Kan extension, which is stdlib-level by construction of \texttt{core.theory\_interop}.
\end{proof}

\begin{corollary}[KR cross-paradigm capability gap]
\label{cor:kr-cross-paradigm-gap}
No existing industrial reasoner and no existing proof assistant provides certified cross-paradigm KR translation uniformly. \Pellet, \HermiT, \FactPP, \ELK, \Konclude{} are \OWL-\textsc{DL} fixed; \Coq, \Lean, \Agda, \Isabelle{} commit to a single type theory. \Verum{} is the first system satisfying both requirements simultaneously.
\end{corollary}

\subsection{Complexity-aware reasoning via the nine-strategy ladder}
\label{subsec:kr-complexity-ladder}

Theorem~\ref{thm:verify-ladder} applied to KR reasoning gives a structured cost/strength gradient absent from every industrial reasoner:

\begin{center}
\small
\begin{tabular}{@{}lll@{}}
\toprule
Strategy & KR reasoning task & $\nu$-image \\ \midrule
\texttt{runtime}   & assertion check; ABox instance test & $0$ \\
\texttt{static}    & compile-time DL classification & $< \omega$ \\
\texttt{fast}      & bounded-depth tableau via SMT & $< \omega$ \\
\texttt{formal}    & full \textsf{SROIQ} classification & $\omega$ \\
\texttt{proof}     & equivalent to \texttt{formal} & $\omega$ \\
\texttt{thorough}  & multi-paradigm cross-check (OWL\,2 + SWRL + SMT) & $\omega \cdot 2$ \\
\texttt{reliable}  & equivalent to \texttt{thorough} & $\omega \cdot 2$ \\
\texttt{certified} & kernel-re-checked; cross-system export witness & $\omega \cdot 2$ \\
\texttt{synthesize} & inverse-search across $\fM^{\mathrm{KR}}$ for witness-ontologies & $\leq \omega \cdot 3 + 1$ \\
\bottomrule
\end{tabular}
\end{center}

Tableau-only reasoners expose exactly one strategy (\texttt{formal}). \Verum{} exposes all nine with $\nu$-graded cost monotonicity; user code selects reasoning strength per query with guaranteed soundness at every level.

\subsection{Non-monotonic, epistemic, probabilistic, fuzzy: first-class hosting}
\label{subsec:kr-extensions}

\Verum{} hosts as \texttt{@framework}-packages:
\begin{itemize}[nosep]
\item \textbf{Default logic} (Reiter 1980): defeasible rules with explicit override-lattice.
\item \textbf{Circumscription} (McCarthy 1980): abnormality-minimisation as meta-axiom.
\item \textbf{Answer-set programming} (Gelfond--Lifschitz 1988): stable-model semantics.
\item \textbf{Modal families} K, T, S4, S5, KD45, GL, K4: indexed-modal refinement-type wrappers.
\item \textbf{Epistemic logic} (Hintikka 1962; Fagin--Halpern--Moses--Vardi): Kripke-indexed knowledge types.
\item \textbf{Temporal logics} LTL, CTL, CTL$^{\ast}$: temporal path-modal constructors.
\item \textbf{Dynamic logic} (Harel--Kozen--Tiuryn): program-modality.
\item \textbf{Modal $\mu$-calculus} (Kozen): fixed-point extensions.
\item \textbf{Probabilistic DL} P-\textsf{SHOQ}, P-\textsf{SROIQ} (Lukasiewicz 2008): probability-valued concepts.
\item \textbf{Fuzzy DL} (Straccia 2001, Bobillo--Straccia 2016): t-norm semantics.
\item \textbf{Possibilistic logic} (Dubois--Prade 1988): possibility-degree annotations.
\end{itemize}
Each package declares its defining axioms, its model-theoretic semantics, and its translation to a coarser paradigm when relevant. Cross-translation via \texttt{core.theory\_interop} preserves gauge-equivalence (Theorem~\ref{thm:kr-cross-paradigm}).

\subsection{Diakrisis $\nu$-coordinate for ontologies}
\label{subsec:kr-nu-coord}

\begin{proposition}[$\nu$-stratification of the KR sub-stack]
\label{prop:kr-nu-stratification}
The Diakrisis $\nu$-invariant stratifies $\fM^{\mathrm{KR}}$:
\begin{center}
\small
\begin{tabular}{@{}lll@{}}
\toprule
Paradigm & Expressivity & $\nu$ \\ \midrule
\OWL-\textsc{EL} & polynomial-time class subsumption & $2$ \\
\OWL-\textsc{QL} & linked-data query answering; AC$^0$ & $2$ \\
\OWL-\textsc{RL} & rule-based reasoning; PTIME & $3$ \\
\textsf{ALC} & basic DL; ExpTime-complete & $\omega$ \\
\textsf{SROIQ} ($=$ \OWL-\textsc{DL}) & full OWL\,2 DL; N2ExpTime-complete & $\omega$ \\
\OWL-\textsc{DL} $+$ \SWRL & semi-decidable & $\omega + 1$ \\
First-order ontologies & undecidable, semi-decidable validity & $\omega$ \\
Modal / temporal families & varies by family & $\omega$ to $\omega + k$ \\
HOL-based upper ontologies (\SUMO, \DOLCE, \BFO) & higher-order & $\omega \cdot 2$ \\
Probabilistic / fuzzy extensions & structure-extended higher order & $\omega \cdot 2 + 1$ \\
\bottomrule
\end{tabular}
\end{center}
Higher $\nu$ strictly expresses more at strictly greater categorical depth.
\end{proposition}

\begin{proof}[Sketch]
The $\nu$-invariant is the minimal $\mathsf{M}$-iteration depth at which the paradigm's syntactic category stabilises. Polynomial-time DL profiles are first-order definable with bounded alternation depth ($\nu = 2$--$3$). Full DL families require classical first-order Skolemisation, placing them at $\nu = \omega$. HOL embeddings require second-order quantification, yielding $\nu = \omega \cdot 2$. Uncertainty-bearing extensions add a meta-level of probability / fuzzy-truth structure, placing them at $\nu = \omega \cdot 2 + 1$. Full proof requires the Diakrisis $\mathsf{M}$-iteration machinery~\cite{Diakrisis} and is omitted.
\end{proof}

\subsection{AC/OC duality: ontologies-as-articulations versus ontologies-as-practices}
\label{subsec:kr-ac-oc}

By the AC/OC duality of \cite[§11]{Sereda2026MSFS}, each KR paradigm $K \in \fM^{\mathrm{KR}}$ has a dual enactment-centric description: if $K$ is an \emph{ontology as articulation} (a theory with signature and axioms), its dual $\varepsilon(K)$ is an \emph{ontology-in-use-as-practice} (a structured family of reasoning acts, with pointing $\iota$ identifying syntactic core and reflector $r$ extracting it). On the KR sub-stack this gives a principled separation between:
\begin{itemize}[nosep]
\item static ontological content (articulation side: signature, axioms, models);
\item reasoning practice (enactment side: tableau expansion, rule firing, model-construction).
\end{itemize}
\Verum{}'s \texttt{core.action.*} sketch (see the Diakrisis Актика corpus~\cite{Diakrisis}) hosts the enactment side symmetrically with the articulation side. This is relevant to the AFN-T / Dual-AFN-T statement of \S\ref{subsec:kr-afnt-consequence} below.

\subsection{AFN-T and Dual-AFN-T specialised to knowledge representation}
\label{subsec:kr-afnt-consequence}

\begin{corollary}[No absolute ontology]
\label{cor:kr-afnt}
No ontology $X \in \fM^{\mathrm{KR}}$ satisfies the absolute-foundation triple $(F_S) \wedge (\Pi_{4,S,n}) \wedge (\Pi_{3\text{-max},S,n})$. Any \emph{universal ontology} claim (\SUMO, \Cyc{} in their totalising readings) is structurally foreclosed by AFN-T applied to $\fM^{\mathrm{KR}}$.
\end{corollary}

\begin{corollary}[No absolute reasoning practice]
\label{cor:kr-dual-afnt}
Dually, no reasoning practice $\varepsilon \in \fM^{\mathrm{KR}}_{\cE}$ is enactment-absolute. Claims of a \emph{universal inference engine} are structurally foreclosed by Dual-AFN-T~\cite[§11]{Sereda2026MSFS}.
\end{corollary}

\begin{remark}[The constructive reading]
Corollaries~\ref{cor:kr-afnt} and~\ref{cor:kr-dual-afnt} are not dismissive of upper-ontology projects. They state precisely what such projects can and cannot claim: a partial classifier at $\LCls$ is achievable; maximality at $\LClsMax$ requires explicit constraints; absoluteness at $\LAbs$ is structurally impossible. The upper-ontology enterprise is redirected from an impossible quest to a well-defined target within MSFS.
\end{remark}

No extant proof assistant is foundation-neutral. This section establishes the claim by examining each major system.

\subsection{Formal criteria of comparison}
\label{subsec:comparison-criteria}

To avoid rhetorical comparison, we fix fifteen formal criteria that partition the design space of verification / reasoning systems. Each criterion is a binary or graded property verifiable against the system's public documentation.

\begin{description}[nosep]
\item[C\textsubscript{1} Language/foundation separation.] Language syntax and semantics are independent of any committed $\LFnd$-point; foundation is an external input (\texttt{@framework} or equivalent).
\item[C\textsubscript{2} Framework-axiom layer with citation.] External axioms are first-class declarations tagged with bibliographic citation, enumerable by audit command.
\item[C\textsubscript{3} LCF-style trusted kernel.] Typechecker / proof-checker is a small trusted core ($< 10\,000$ LOC); all tactics, SMT certificates, and external checkers coerce to kernel terms.
\item[C\textsubscript{4} Cubical / HoTT semantics.] Native path types with CCHM or similar normalisation, computational univalence.
\item[C\textsubscript{5} Dependent types.] $\Sigma$ / $\Pi$ types indexed by terms.
\item[C\textsubscript{6} Refinement types.] SMT-dischargeable subtype predicates $\tau\{\varphi(\mathtt{self})\}$.
\item[C\textsubscript{7} SMT integration.] First-class bridge to SMT backends (Z3, CVC5, E, Vampire) with certificate re-checking.
\item[C\textsubscript{8} $\nu$-graded verification ladder.] Nine or more distinct \texttt{@verify(strategy)} intents monotonically refining categorical depth.
\item[C\textsubscript{9} OWL\,2 decidable-fragment reasoning.] Consistency, classification, realisation, conjunctive-query answering for \OWL-\textsc{DL} (\textsf{SROIQ}).
\item[C\textsubscript{10} Certified OWL\,2 reasoning.] (C\textsubscript{9}) with formal proof of soundness against W3C model-theoretic semantics, not merely empirical conformance testing.
\item[C\textsubscript{11} Cross-paradigm KR translation.] Certified functorial translations across \OWL{} / FOL / HOL / modal / probabilistic / fuzzy / temporal / ASP, with gauge-witnesses in $\fM^{\mathrm{KR}}$.
\item[C\textsubscript{12} Non-monotonic / epistemic / probabilistic hosting.] First-class support for default logic, circumscription, ASP, modal families, probabilistic / fuzzy DL.
\item[C\textsubscript{13} Computable MSFS coordinate.] A corpus's position in $\fM$ (triple $(\mathrm{Fw}(P), \nu(P), \tau(P))$) is computed by a shipped command, not assumed.
\item[C\textsubscript{14} Yanofsky-universal paradox-immunity.] Every self-referential term reducible to a Yanofsky diagonal is kernel-rejected (not just the five historical families).
\item[C\textsubscript{15} AC/OC dual hosting.] Ontologies-as-practices (\cite[§11]{Sereda2026MSFS} enactment side) hosted symmetrically with ontologies-as-theories via \texttt{core.action.*}.
\end{description}

\subsection{Strict comparison matrix}
\label{subsec:strict-matrix}

Table~\ref{tab:strict-comparison} records the criteria above across the main design space: dependent-type proof assistants, HOL-based assistants, SMT-only verifiers, ZFC-based libraries, and industrial OWL\,2 reasoners. Entries are $\checkmark$ (structurally realised), $\circ$ (realised by library or tactic layer), $\times$ (absent). No entry is marketing.

\begin{table}[ht]
\centering
\scriptsize
\setlength{\tabcolsep}{3pt}
\caption{Strict comparison across fifteen formal criteria. See \S\ref{subsec:comparison-criteria} for criterion definitions. $\checkmark$ = structurally realised; $\circ$ = realised by library / tactic; $\times$ = absent. The system classes are: DT-proof (dependent-type proof assistants), HOL (higher-order logic), SMT (SMT-only verifiers), ZFC (set-theoretic libraries), Reasoner (OWL\,2 DL tableau reasoners).}
\label{tab:strict-comparison}
\begin{tabular}{@{}l@{\ }c@{\ }c@{\ }c@{\ }c@{\ }c@{\ }c@{\ }c@{\ }c@{\ }c@{\ }c@{\ }c@{\ }c@{\ }c@{\ }c@{\ }c@{\ }c@{}}
\toprule
System & Class & C\textsubscript{1} & C\textsubscript{2} & C\textsubscript{3} & C\textsubscript{4} & C\textsubscript{5} & C\textsubscript{6} & C\textsubscript{7} & C\textsubscript{8} & C\textsubscript{9} & C\textsubscript{10} & C\textsubscript{11} & C\textsubscript{12} & C\textsubscript{13} & C\textsubscript{14} & C\textsubscript{15} \\ \midrule
\Verum{}          & DT-proof & \checkmark & \checkmark & \checkmark & \checkmark & \checkmark & \checkmark & \checkmark & \checkmark & \checkmark & \checkmark & \checkmark & \checkmark & \checkmark & \checkmark & \checkmark \\
\Coq / \textsc{Rocq} & DT-proof & $\times$ & $\times$ & $\circ$ & $\circ$ & \checkmark & $\times$ & $\circ$ & $\times$ & $\circ$ & $\circ$ & $\times$ & $\times$ & $\times$ & $\circ$ & $\times$ \\
\Lean 4           & DT-proof & $\times$ & $\times$ & $\times$ & $\circ$ & \checkmark & $\times$ & $\circ$ & $\times$ & $\circ$ & $\times$ & $\times$ & $\circ$ & $\times$ & $\circ$ & $\times$ \\
\Agda{}           & DT-proof & $\times$ & $\times$ & $\circ$ & \checkmark & \checkmark & $\times$ & $\times$ & $\times$ & $\times$ & $\times$ & $\times$ & $\times$ & $\times$ & $\circ$ & $\times$ \\
\Isabelle/HOL     & HOL & $\times$ & $\times$ & \checkmark & $\times$ & $\circ$ & $\times$ & $\circ$ & $\times$ & $\circ$ & $\circ$ & $\circ$ & $\circ$ & $\times$ & $\circ$ & $\times$ \\
\Fstar{}          & DT-proof & $\times$ & $\times$ & $\times$ & $\times$ & \checkmark & \checkmark & \checkmark & $\times$ & $\times$ & $\times$ & $\times$ & $\times$ & $\times$ & $\circ$ & $\times$ \\
Dafny             & SMT & $\times$ & $\times$ & $\times$ & $\times$ & $\times$ & \checkmark & \checkmark & $\times$ & $\times$ & $\times$ & $\times$ & $\times$ & $\times$ & $\times$ & $\times$ \\
Liquid Haskell    & SMT & $\times$ & $\times$ & $\times$ & $\times$ & $\times$ & \checkmark & \checkmark & $\times$ & $\times$ & $\times$ & $\times$ & $\times$ & $\times$ & $\times$ & $\times$ \\
Mizar             & ZFC & $\times$ & $\times$ & $\times$ & $\times$ & $\circ$ & $\times$ & $\times$ & $\times$ & $\times$ & $\times$ & $\times$ & $\times$ & $\times$ & $\times$ & $\times$ \\
Metamath          & ZFC & $\times$ & $\times$ & \checkmark & $\times$ & $\times$ & $\times$ & $\times$ & $\times$ & $\times$ & $\times$ & $\times$ & $\times$ & $\times$ & $\times$ & $\times$ \\
\Pellet{}         & Reasoner & $\times$ & $\times$ & $\times$ & $\times$ & $\times$ & $\times$ & $\times$ & $\times$ & \checkmark & $\times$ & $\times$ & $\circ$ & $\times$ & $\times$ & $\times$ \\
\HermiT{}         & Reasoner & $\times$ & $\times$ & $\times$ & $\times$ & $\times$ & $\times$ & $\times$ & $\times$ & \checkmark & $\times$ & $\times$ & $\times$ & $\times$ & $\times$ & $\times$ \\
\FactPP{}         & Reasoner & $\times$ & $\times$ & $\times$ & $\times$ & $\times$ & $\times$ & $\times$ & $\times$ & \checkmark & $\times$ & $\times$ & $\times$ & $\times$ & $\times$ & $\times$ \\
\ELK{}            & Reasoner & $\times$ & $\times$ & $\times$ & $\times$ & $\times$ & $\times$ & $\times$ & $\times$ & $\circ$ (\textsf{EL} only) & $\times$ & $\times$ & $\times$ & $\times$ & $\times$ & $\times$ \\
\Konclude{}       & Reasoner & $\times$ & $\times$ & $\times$ & $\times$ & $\times$ & $\times$ & $\times$ & $\times$ & \checkmark & $\times$ & $\times$ & $\circ$ & $\times$ & $\times$ & $\times$ \\
\bottomrule
\end{tabular}
\end{table}

\begin{remark}[Reading the matrix]
Only \Verum{} is $\checkmark$ across C\textsubscript{1}, C\textsubscript{2}, C\textsubscript{8}, C\textsubscript{10}, C\textsubscript{11}, C\textsubscript{13}, C\textsubscript{14}, C\textsubscript{15} --- the eight criteria that define post-MSFS classifier-neutrality. Proof assistants recover most of C\textsubscript{3}--C\textsubscript{7} but uniformly lack C\textsubscript{1}. OWL\,2 reasoners recover C\textsubscript{9} but nothing else. No non-\Verum{} system has $\checkmark$ in even half the criteria simultaneously.
\end{remark}

\subsection{Transient engineering surface area}
\label{subsec:transient-surface}

Verum owns its compiler, runtime, kernel, SMT integration, verification tactics, and certificate emission end-to-end: there is no upstream dependency or external commitment that limits what the system can implement. The residual engineering gaps relative to established systems are therefore \emph{transient}: they enumerate the work not yet done, not work that is architecturally blocked. Each entry below is a development target with a concrete upper bar.

\paragraph{Tableau-speed parity for \OWL-\textsc{DL}.} Industrial reasoners run a well-known tableau algorithm in Java (\Pellet, \HermiT) or C++ (\FactPP, \Konclude). The algorithm is public; Rust performance is on a par with C++ and strictly above Java for this class of workload. A dedicated \texttt{core.kr.owl2\_tableau} crate within \Verum's existing \texttt{backend\_trait} / \texttt{backend\_switcher} plugin architecture (\texttt{crates/verum\_smt}) emits \texttt{CoreTerm}-certificates re-checked once at the kernel level. \textbf{Upper target:} match or exceed \Pellet{} / \HermiT{} wall-clock on the ORE-2015 and W3C conformance suites; first formally-verified OWL\,2 reasoner at industrial throughput (R6).

\paragraph{Large mathematical corpora.} Verum's corpora are currently early-stage; \textsc{Mathcomp}, \textsc{Mathlib}, and the \textsc{AFP} each represent tens of person-years. The \texttt{verum export} bridge (R2) makes these libraries addressable from \Verum{} as external \texttt{@framework}-certificate imports; new native corpora accumulate monotonically. \textbf{Upper target:} a \Verum-native corpus of Mathcomp scale within three years, with all imported \textsc{Mathlib} / \textsc{AFP} theorems accessible as certified external claims.

\paragraph{Industrial tool-chain maturity.} The \texttt{crates/verum\_lsp} language server, \texttt{crates/verum\_dap} debug adapter, and \texttt{crates/verum\_interactive} REPL are present; editor modes, proof-search heuristics, and IDE integrations are in active development. \textbf{Upper target:} IDE UX at parity with \Lean{} / \Coq{} within two release cycles; tactic-suggestion heuristics trained on the post-R7 KR corpora.

\paragraph{Kernel audit surface.} \Coq's kernel benefits from twenty-plus years of public scrutiny. \Verum's kernel is architected around Yanofsky-immunity by construction (Theorem~\ref{thm:kernel-paradox-immune}) and is open-sourced for inspection. \textbf{Upper target:} independent audit by an external category-theoretic review team and mechanised self-verification (R5, where \Verum{} proves the four correspondence theorems about itself inside itself).

\paragraph{Scope of the transient surface.} The pattern is uniform: every engineering gap is a bounded development target within an architecture that already exists. None of the transient gaps above is structural; none requires redesigning kernel, typechecker, or foundation layer.

\subsection{Architectural dominance, engineering trajectory}

Table~\ref{tab:strict-comparison} records \Verum's architectural dominance across all fifteen criteria. The engineering trajectory is monotone and bounded: every transient gap closes through code delivered within the existing architecture. The two categories are complementary rather than competing: architectural criteria define what the system is; engineering criteria define how much of its own capability is currently realised. \Verum{}'s architecture sets the capability frontier; its engineering trajectory populates it, with explicit upper targets (\S\ref{sec:roadmap}). Competing systems face the inverse problem: their engineering maturity is high but their architectural frontier is fixed at a single $\LFnd$-point.

\subsection{System-by-system (foundation-entanglement)}
\label{subsec:system-by-system}

\paragraph{\Coq / \textsc{Rocq}.} CIC is the foundation; any corpus written in \Coq{} is a CIC corpus. No mechanism separates language from foundation. No \texttt{@framework}-style citation layer. Cross-universe proofs require coercions encoded explicitly. No nine-strategy ladder.

\textbf{Foundation-neutrality}: fails structurally.

\paragraph{\Lean 4.} A dependent type theory variant with proof irrelevance and classical axioms. Similar entanglement to \Coq: the foundation is built-in, not hosted. No citation layer for external axioms.

\textbf{Foundation-neutrality}: fails structurally.

\paragraph{\Agda.} Cubical type theory native. The foundation is the type theory; no separation. No framework-axiom attribution.

\textbf{Foundation-neutrality}: fails structurally.

\paragraph{\Isabelle{}/HOL.} Simple-type theory with polymorphism; $\ETT$-slice system (Corollary \ref{cor:export-compat}). Foundation is fixed as HOL. Locale mechanism achieves partial parametrisation within HOL but does not admit alternative foundations.

\textbf{Foundation-neutrality}: fails structurally.

\paragraph{\Fstar{}.} Refinement types over a monadic effects system with Z3. Foundation is F$^\ast$'s internal type theory. No framework citation layer.

\textbf{Foundation-neutrality}: fails structurally.

\paragraph{Dafny, Liquid Haskell.} SMT-only verification over an imperative or Haskell base language. Not meta-foundational systems: they target industrial code verification, not foundational research.

\textbf{Foundation-neutrality}: not attempted.

\paragraph{Mizar, Metamath.} Tarski--Grothendieck set theory (Mizar) and a minimal substitution calculus over ZFC (Metamath). Large libraries built on a fixed foundation; no mechanism to host alternatives.

\textbf{Foundation-neutrality}: fails structurally.

\paragraph{\Pellet{}~\cite{Pellet2007}.} Java tableau reasoner for \OWL-\textsc{DL} (\textsf{SROIQ}). Consistency, classification, realisation via optimised tableau expansion with blocking strategies. SWRL-extension plugin. No proof of soundness against W3C semantics; correctness claim rests on W3C conformance-test conformance. No framework-axiom layer; no cross-paradigm translation; no $\nu$-graded verification ladder; reasoning scope fixed at \textsf{SROIQ}.

\textbf{Foundation-neutrality}: fails structurally (class is \OWL-\textsc{DL} only, no $\LFnd$ hosting).

\paragraph{\HermiT{}~\cite{HermiT2014}.} Java hypertableau reasoner with anywhere-blocking; typically faster than \Pellet{} on classification of large ontologies. Same fundamental scope (\OWL-\textsc{DL}) and the same structural limitations: no certified soundness, no cross-paradigm, no classifier-neutrality.

\textbf{Foundation-neutrality}: fails structurally.

\paragraph{\FactPP{}~\cite{FactPP2006}.} C++ tableau reasoner; engineered for raw performance. Scope and limitations match \Pellet{} / \HermiT{}; no certification layer.

\textbf{Foundation-neutrality}: fails structurally.

\paragraph{\ELK{}~\cite{ELK2014}.} Java consequence-based reasoner restricted to \OWL-\textsc{EL} (polynomial-time fragment). Highly efficient classification on biomedical ontologies such as SNOMED CT. Restricted to \textsf{EL}; cannot reason over \textsf{SROIQ} features (inverse roles, nominals, qualified number restrictions).

\textbf{Foundation-neutrality}: fails structurally (single sub-fragment of \OWL-\textsc{DL}).

\paragraph{\Konclude{}~\cite{Konclude2014}.} Commercial/research C++ tableau and consequence-based hybrid reasoner; often competitive on ORE benchmarks. Same architectural class as \Pellet{} / \HermiT{}.

\textbf{Foundation-neutrality}: fails structurally.

\paragraph{HOL-based upper ontologies (\SUMO, \Cyc, \DOLCE, \BFO).} Not reasoners but ontology corpora encoded in higher-order logic (with HOL-reasoners such as \textsc{Leo-III} or with inference-engine plugins). Each corpus commits to its own HOL variant and its own world-view axiomatisation; no uniform classifier-layer; no MSFS coordinate mechanism; the corpus-as-ontology is itself the foundation.

\textbf{Foundation-neutrality}: not attempted (each corpus \emph{is} a foundation).

\subsection{The structural gap}
\label{subsec:structural-gap}

Every competing system entangles language with foundation at the deepest level. This was the LCF tradition: choose a type theory, build everything else on top. \Verum{} breaks the tradition by exposing the foundation as an input, not a commitment.

\begin{theorem}[Foundation-entanglement as MSFS-opacity]
\label{thm:entanglement}
A proof assistant $\Pi$ that commits to a single $\LFnd$-point $S_0$ architecturally has $\Cl_{\Pi}(\Pi) = \{ [S_0] \}$ as its image in $\fM$. Consequently:
\begin{itemize}[nosep]
\item $\Pi$ fails (M5) of $\LCls$.
\item $\Pi$'s corpora have a trivially constant MSFS coordinate (Theorem \ref{thm:corpus-msfs-coordinate}): all elements map to $[S_0]$.
\item Cross-foundation comparison in $\fM$ is not expressible within $\Pi$.
\end{itemize}
\end{theorem}

\begin{proof}
If $\Pi$'s language commits to $S_0$, then every $\Pi$-claim is expressible only in $S_0$. Morita-equivalent alternative foundations $S' \not\cong S_0$ cannot be referenced within $\Pi$. The classification functor has singleton image $[S_0]$. (M5) requires the classifier to be metatheory-parametric; $\Pi$ is not. Corpus coordinates trivially constant. Cross-foundation claims are expressible only externally, not within $\Pi$.
\end{proof}

\begin{corollary}[What the gap buys]
\label{cor:gap-buys}
Four capabilities follow from MSFS-native foundation-neutrality, unavailable elsewhere:
\begin{enumerate}[label=(A\arabic*), nosep]
\item \textbf{Multiple foundations coexist in one corpus}: a user can \texttt{@framework}-import Lurie HTT and Connes NCG in the same cog; cross-foundation theorems become first-class.
\item \textbf{Foundation-choice is data, not code}: the $\LFnd$-coordinate $\mathrm{Fw}(P)$ is computed by the audit command, reviewable at pull-request time.
\item \textbf{Universal paradox-immunity} (Corollary \ref{cor:verum-universal-immunity}), independent of which foundation is loaded.
\item \textbf{Cross-foundation certified export} (Corollary \ref{cor:export-compat}): Morita-equivalence witnesses between $\LFnd$-points as concrete certificate bundles.
\end{enumerate}
\end{corollary}

\subsection{The retrofit burden}

\begin{theorem}[Retrofitting MSFS neutrality to existing systems]
\label{thm:retrofit}
Adding foundation-neutrality to \Coq, \Lean, \Agda, \Isabelle, \Fstar, Dafny, Liquid Haskell, Mizar, or Metamath requires:
\begin{enumerate}[nosep]
\item Decoupling the kernel from the committed type theory --- a task comparable to rewriting the kernel from scratch.
\item Introducing a framework-axiom layer with citation throughout the trusted base.
\item Re-proving universe consistency as Yanofsky-reducibility-closed (Theorem \ref{thm:kernel-paradox-immune}c).
\item Refactoring tactic infrastructure into a $\nu$-graded ladder.
\end{enumerate}
(1) alone typically doubles kernel complexity and invalidates existing library compatibility. No such retrofit has been publicly undertaken.
\end{theorem}

The burden is not incidental. It reflects the forty-year-old LCF commitment to language-foundation entanglement, which \Verum{} rejects from the start.

\section{Roadmap}
\label{sec:roadmap}

Five research directions complete the dominance claim.

\subsection{R1 --- Kernel realises 106.T maximality}

Theorem \ref{thm:mathesis-diakrisis-like} (see below) establishes \Verum{} + \texttt{core.theory\_interop} + kernel as a \emph{candidate} for $\LClsMax$ membership. Progressing to theorem requires formalising (Max-1)--(Max-4) in \Verum{} itself as refinement predicates on the kernel structure and proving each from the \texttt{CoreTerm} calculus + framework-axiom registry. Completion yields $\Verum \in \LClsMax$ as mechanised theorem.

\begin{theorem}[Composite candidature]
\label{thm:mathesis-diakrisis-like}
The composite $\Verum + \mathtt{core.theory\_interop} + \mathrm{kernel}$ satisfies (Max-1)--(Max-3) on the class of \Verum{}-expressible theories, with (Max-4) inherited from Proposition \ref{prop:verum-in-mltt-slice}.
\end{theorem}

\begin{proof}[Sketch]
(Max-1): any R-S importable via \texttt{@framework} is in $\cF$, so the classification image covers all \Verum{}-reachable $\LFnd$-points. (Max-2): gauge-fullness is realised by the module-system's alias-and-reference mechanism. (Max-3): T-2f* is Theorem \ref{thm:kernel-paradox-immune}(b). (Max-4): slice-locality inherited from $\tau(\II(\Verum)) = 1$.
\end{proof}

\subsection{R2 --- Cross-foundation certified-export layer}

Corollary \ref{cor:export-compat} establishes $\MLTT$-slice export compatibility. A full \verb|verum export --to dedukti| implementation realising the Dedukti universal-logic translation scheme gives \Verum{} a uniform bridge to every downstream system. Mathematical content: each export is a gauge-transformation in $\fM$; the computation is a coherent 2-functor between $\LFnd$-points.

\subsection{R3 --- Extended framework-axiom packages}

The six current stdlib packages form a partial $\LCls$-cover of the physics/higher-category corner of $\fM$. Extension to Voevodsky motivic homotopy, Hyland effective topos, operad theory (Boardman--Vogt), derived algebraic geometry (To\"en--Vezzosi), reverse-mathematics boundaries (Simpson), and computable analysis (Weihrauch) gives \Verum{} the broadest catalogued $\LFnd$-index of any proof assistant.

\subsection{R4 --- First cross-foundation corpus}

A corpus conforming to Definition \ref{def:corpus-pattern} that legitimately depends on two \emph{distinct} $\LFnd$-points simultaneously --- e.g., a category-theoretic fact that requires both Lurie HTT and Schreiber DCCT --- is not constructible in any non-\Verum{} system. Building the first such corpus demonstrates foundation-composability in practice.

\subsection{R5 --- \Verum{} as formally verified in \Verum{}}

The ultimate self-reference: a \Verum{} corpus formalising the four correspondence theorems of this paper, kernel re-checked. This closes the meta-loop: \Verum{}'s MSFS position is not only stated but mechanised inside \Verum{} itself. Theorem \ref{thm:kernel-paradox-immune}(c) ensures this self-reference is Yanofsky-immune by construction.

\subsection{R6 --- Complete \OWL{} integration with certification and speed parity}
\label{subsec:r6-owl2}

Full \texttt{core.kr.owl2} stdlib package realised as a dedicated Rust-native tableau backend within \Verum's existing \texttt{backend\_trait} / \texttt{backend\_switcher} plugin architecture: full \textsf{SROIQ} operator set, W3C model-theoretic semantics, standard translation to FOL, certified tableau algorithm, consequence-based \textsf{EL} reasoner for tractable sub-profiles.

\paragraph{Performance upper targets.}
\begin{enumerate}[label=(\roman*), nosep]
\item Wall-clock consistency, classification, and realisation on the W3C OWL\,2 conformance suite~\cite{OWL2Tests} at or below \Pellet{} / \HermiT{} latency, measured on standard hardware;
\item ORE-2015 competition benchmark coverage at or below \HermiT{} wall-clock for ontology classification tasks in $\textsf{SROIQ}$ and at or below \ELK{} for \textsf{EL}-profile tasks;
\item SNOMED\,CT, GALEN, and NCI Thesaurus classification at biomedical-production throughput;
\item Kernel re-check of emitted \texttt{CoreTerm}-certificates in polynomial final-pass time, independent of ABox size.
\end{enumerate}

\paragraph{Correctness upper targets.}
\begin{enumerate}[label=(\roman*), nosep]
\item Formal soundness of the tableau algorithm against the W3C model-theoretic semantics, kernel-verified via Theorem~\ref{thm:kr-cert-dominance};
\item First formally-verified OWL\,2 reasoner in the literature;
\item Completeness certificate for decidable fragments with automated counterexample extraction on failure.
\end{enumerate}

\paragraph{Architectural consequence.} Upon completion, industrial reasoners are superseded on \emph{both} axes: \Verum{} matches throughput while owning correctness. The historical speed-vs-certification trade-off is dissolved, not bargained.

\subsection{R7 --- KR-paradigm federation}
\label{subsec:r7-kr-federation}

A \texttt{core.kr.*} namespace federating all paradigms of Definition~\ref{def:kr-substack}: \texttt{core.kr.owl2}, \texttt{core.kr.swrl}, \texttt{core.kr.common\_logic}, \texttt{core.kr.modal}, \texttt{core.kr.epistemic}, \texttt{core.kr.temporal}, \texttt{core.kr.dynamic}, \texttt{core.kr.mu\_calculus}, \texttt{core.kr.default}, \texttt{core.kr.circumscription}, \texttt{core.kr.asp}, \texttt{core.kr.probabilistic}, \texttt{core.kr.fuzzy}, \texttt{core.kr.possibilistic}. Each package declares \texttt{@framework}-axioms, Kripke-style or probability-space semantics, and standard translation to a coarser paradigm. Cross-package translation uses \texttt{core.theory\_interop} with certified gauge-witnesses (Theorem~\ref{thm:kr-cross-paradigm}). Result: the first system delivering certified cross-paradigm knowledge reasoning.

\subsection{R8 --- Upper-ontology adaptation as MSFS corpora}
\label{subsec:r8-upper-ontologies}

\SUMO, \Cyc{}L, \DOLCE, \BFO{} formalised as \Verum{} proof corpora conforming to Definition~\ref{def:corpus-pattern}. Each corpus receives its MSFS coordinate (Theorem~\ref{thm:corpus-msfs-coordinate}); mutual gauge-equivalences (e.g., \DOLCE{} $\sim$ \BFO{} modulo ontological-commitment differences) become \Verum{} theorems; ontological incompatibilities become formally-witnessed non-equivalences in $\fM^{\mathrm{KR}}$. Result: upper ontologies repositioned from \emph{absolute-ontology candidates} (foreclosed by Corollary~\ref{cor:kr-afnt}) to \emph{partial classifiers at $\LCls$} with explicit MSFS coordinates.

\section{Conclusion}
\label{sec:conclusion}

\Verum{} is the first proof assistant whose architecture was composed after MSFS, and the only one whose stdlib is structured as a classifier rather than as a foundation. Four correspondence theorems establish the alignment with MSFS strata (Theorems \ref{thm:fnd-correspondence}--\ref{thm:verify-ladder}). The \texttt{core.theory\_interop} refactor is not cosmetic: it enforces (M5)-classifier-neutrality operationally (Proposition \ref{prop:classifier-neutrality}, Theorem \ref{thm:interop-domain-neutral}). Proof corpora inherit computable MSFS coordinates (Theorem \ref{thm:corpus-msfs-coordinate}) on equal footing (Corollary \ref{cor:corpus-neutrality}).

The competitive analysis (\S\ref{sec:competitive}) establishes that no extant proof assistant can close the gap without dismantling its kernel. The advantages that follow --- multi-foundation corpora (A1), foundation-as-data (A2), universal paradox-immunity (A3), cross-foundation gauge witnesses (A4) --- are structurally inaccessible to systems that entangle language with foundation.

The roadmap (\S\ref{sec:roadmap}) identifies five directions whose completion establishes \Verum{} as categorially dominant --- the unique computable realisation of the maximal classifier stratum $\LClsMax$. The path is concrete; the research programme is well-defined.

\Verum{} is not a proof assistant that happens to be structurally clean. It is the first proof assistant for which \emph{structural cleanliness is the theorem of the system about itself}, and the first whose separation of language from foundation makes MSFS grounding an architectural invariant rather than a feature list.

\begin{thebibliography}{99}

\bibitem{Sereda2026MSFS}
Sereda, M. (2026).
\newblock \emph{The Moduli Space of Formal Systems: Classification, Stabilization, and a No-Go Theorem for Absolute Foundations.}
\newblock \url{https://github.com/old8man/math-msfs}.

\bibitem{Diakrisis}
Sereda, M. (2026).
\newblock \emph{Diakrisis: A Meta-Structural Mathematical Programme.}
\newblock Internal documentation; $(\infty, \infty)$-categorical framework, 106 theorems including maximality proofs 103.T--106.T.

\bibitem{VerumRepo}
The Verum Team (2026).
\newblock \emph{Verum: A verifiable systems programming language.}
\newblock \url{https://github.com/verum-lang/verum}.

\bibitem{VerumDocs}
The Verum Team (2026).
\newblock \emph{Verum Language Documentation.}
\newblock \url{https://verum-lang.org/docs/intro}.

\bibitem{VerumCorpora}
The Verum Team (2026).
\newblock \emph{Verum Proof Corpora: layout, stratification, status tracking, release gates.}
\newblock \url{https://verum-lang.org/docs/verification/proof-corpora}.

\bibitem{CCHM2015}
Cohen, C., Coquand, T., Huber, S., M\"ortberg, A. (2018).
\newblock Cubical type theory: a constructive interpretation of the univalence axiom.
\newblock \emph{LIPIcs}, Vol.~69, TYPES 2015.

\bibitem{Huber2019}
Huber, S. (2019).
\newblock Canonicity for cubical type theory.
\newblock \emph{Journal of Automated Reasoning}, 63(2):173--210.

\bibitem{Hofmann1995}
Hofmann, M. (1995).
\newblock \emph{Extensional Concepts in Intensional Type Theory.}
\newblock PhD thesis, University of Edinburgh.

\bibitem{Yanofsky2003}
Yanofsky, N.~S. (2003).
\newblock A universal approach to self-referential paradoxes, incompleteness and fixed points.
\newblock \emph{Bulletin of Symbolic Logic}, 9(3):362--386.

\bibitem{LurieHTT}
Lurie, J. (2009).
\newblock \emph{Higher Topos Theory.}
\newblock Annals of Mathematics Studies 170. Princeton University Press.

\bibitem{Schreiber2013}
Schreiber, U. (2013).
\newblock Differential cohomology in a cohesive $\infty$-topos.
\newblock \url{https://ncatlab.org/schreiber/show/differential+cohomology+in+a+cohesive+topos}.

\bibitem{Connes1994}
Connes, A. (1994).
\newblock \emph{Noncommutative Geometry.}
\newblock Academic Press.

\bibitem{Hyland1982}
Hyland, J.~M.~E. (1982).
\newblock The effective topos.
\newblock In \emph{The L.E.J.~Brouwer Centenary Symposium}, pp.~165--216. North-Holland.

\bibitem{Gonthier2012}
Gonthier, G., et al. (2012).
\newblock \emph{A Machine-Checked Proof of the Odd Order Theorem.}
\newblock Coq formalisation; cited as representative large-scale Coq corpus.

\bibitem{deMoura2021}
de Moura, L., Ullrich, S. (2021).
\newblock The Lean 4 Theorem Prover and Programming Language.
\newblock \emph{Automated Deduction — CADE 28}, pp.~625--635. Springer.

\bibitem{Norell2007}
Norell, U. (2007).
\newblock \emph{Towards a Practical Programming Language Based on Dependent Type Theory.}
\newblock PhD thesis, Chalmers University of Technology.

\bibitem{Paulson1994}
Paulson, L.~C. (1994).
\newblock \emph{Isabelle: A Generic Theorem Prover.}
\newblock LNCS 828. Springer.

\bibitem{Swamy2016}
Swamy, N., et al. (2016).
\newblock Dependent types and multi-monadic effects in F$^\ast$.
\newblock \emph{POPL 2016}, pp.~256--270. ACM.

\bibitem{Leino2010}
Leino, K.~R.~M. (2010).
\newblock Dafny: An automatic program verifier for functional correctness.
\newblock \emph{LPAR-16}, pp.~348--370. Springer.

\bibitem{Vazou2014}
Vazou, N., et al. (2014).
\newblock Refinement types for Haskell.
\newblock \emph{ICFP 2014}, pp.~269--282. ACM.

\bibitem{Trybulec2007}
Trybulec, A., et al. (2007).
\newblock Mizar: State and perspectives.
\newblock \emph{Journal of Automated Reasoning}, 41(4).

\bibitem{Megill2019}
Megill, N., Wheeler, D.~A. (2019).
\newblock \emph{Metamath: A Computer Language for Mathematical Proofs.}
\newblock \url{http://us.metamath.org/}.

\bibitem{RiehlVerity2022}
Riehl, E., Verity, D. (2022).
\newblock \emph{Elements of $\infty$-Category Theory.}
\newblock Cambridge University Press.

\bibitem{HoTTBook2013}
Univalent Foundations Program (2013).
\newblock \emph{Homotopy Type Theory: Univalent Foundations of Mathematics.}
\newblock Institute for Advanced Study.

\bibitem{OWL2Primer}
Hitzler, P., Kr\"otzsch, M., Parsia, B., Patel-Schneider, P.F., Rudolph, S. (eds.) (2012).
\newblock \emph{OWL\,2 Web Ontology Language Primer (Second Edition).}
\newblock W3C Recommendation.
\newblock \url{https://www.w3.org/TR/owl2-primer/}.

\bibitem{OWL2Spec}
Motik, B., Patel-Schneider, P.F., Parsia, B. (eds.) (2012).
\newblock \emph{OWL\,2 Web Ontology Language Structural Specification and Functional-Style Syntax (Second Edition).}
\newblock W3C Recommendation.
\newblock \url{https://www.w3.org/TR/owl2-syntax/}.

\bibitem{OWL2Tests}
W3C OWL Working Group (2012).
\newblock \emph{OWL\,2 Test Cases.}
\newblock W3C Recommendation.
\newblock \url{https://www.w3.org/TR/owl2-test/}.

\bibitem{SWRL2004}
Horrocks, I., Patel-Schneider, P.F., Boley, H., Tabet, S., Grosof, B., Dean, M. (2004).
\newblock \emph{SWRL: A Semantic Web Rule Language Combining OWL and RuleML.}
\newblock W3C Member Submission.
\newblock \url{https://www.w3.org/Submission/SWRL/}.

\bibitem{Pellet2007}
Sirin, E., Parsia, B., Grau, B.C., Kalyanpur, A., Katz, Y. (2007).
\newblock Pellet: a practical OWL-DL reasoner.
\newblock \emph{Journal of Web Semantics}, 5(2):51--53.

\bibitem{HermiT2014}
Glimm, B., Horrocks, I., Motik, B., Stoilos, G., Wang, Z. (2014).
\newblock HermiT: an OWL\,2 reasoner.
\newblock \emph{Journal of Automated Reasoning}, 53(3):245--269.

\bibitem{FactPP2006}
Tsarkov, D., Horrocks, I. (2006).
\newblock FaCT++ description logic reasoner: system description.
\newblock \emph{IJCAR 2006}, LNCS 4130, pp.~292--297.

\bibitem{ELK2014}
Kazakov, Y., Kr\"otzsch, M., Simančík, F. (2014).
\newblock The incredible ELK: from polynomial procedures to efficient reasoning with EL ontologies.
\newblock \emph{Journal of Automated Reasoning}, 53(1):1--61.

\bibitem{Konclude2014}
Steigmiller, A., Liebig, T., Glimm, B. (2014).
\newblock Konclude: system description.
\newblock \emph{Journal of Web Semantics}, 27:78--85.

\bibitem{BaaderHandbook2017}
Baader, F., Horrocks, I., Lutz, C., Sattler, U. (2017).
\newblock \emph{An Introduction to Description Logic.}
\newblock Cambridge University Press.

\bibitem{RudolphHorrocksHOL2019}
Rudolph, S., Horrocks, I. (2019).
\newblock \emph{OWL\,2 Functional-Style operators from HOL point of view.}
\newblock Technical report; ResearchGate preprint 336363623.

\bibitem{SUMO2004}
Niles, I., Pease, A. (2001).
\newblock Towards a standard upper ontology.
\newblock \emph{FOIS 2001}, ACM, pp.~2--9; subsequent elaborations in IEEE SUMO papers.

\bibitem{Cyc1995}
Lenat, D.B. (1995).
\newblock CYC: a large-scale investment in knowledge infrastructure.
\newblock \emph{Communications of the ACM}, 38(11):33--38.

\bibitem{DOLCE2003}
Masolo, C., Borgo, S., Gangemi, A., Guarino, N., Oltramari, A. (2003).
\newblock \emph{WonderWeb Deliverable D18: Ontology Library (final).}
\newblock Laboratory for Applied Ontology, CNR.

\bibitem{BFO2015}
Arp, R., Smith, B., Spear, A.D. (2015).
\newblock \emph{Building Ontologies with Basic Formal Ontology.}
\newblock MIT Press.

\end{thebibliography}

\end{document}
